\summarychapter{10}{Splitting Important Hairs}{The Doctrine of the Trinity}

\begin{inbrief}
The most distinctive thing Christians say about God is that he is three persons. The fourth century
gave this belief its form. The chapter recounts the Arian dispute, the Council of Nicaea and its creed,
the half-century of struggle in which Athanasius defended it, and the analogies by which the church
learned to speak of three in one.
\question{What did the church mean by the Trinity? What is the orthodox understanding of the triune God?}
\end{inbrief}

No other major religion worships a three-in-one God: Muslims and Jews find the doctrine offensive,
Unitarians and Jehovah's Witnesses deplorable, and most Christians admit it is a mystery. In worship
believers spoke of Father, Son and Spirit; thinking about God's nature, they said he is one. In the
fourth century the question was everywhere: in Constantinople, a bishop complained, the money-changer,
the baker and the bath attendant all argued whether the Son is begotten or created, equal or less.

\ssection{Making Sense of Mystery}
Constantine hoped the church would give the weary empire new life, but only a united church could do
so, and he was disturbed that Christians who had just survived persecution now demanded that the state
suppress their opponents. The worst quarrel was at Alexandria. About 318 Arius, pastor of the
influential Baucalis church, taught that the Word who became flesh in Jesus was a lesser god of a
different nature from the Father, neither eternal nor omnipotent: the first and greatest of created
beings, but created. ``The Son has a beginning, but God is without beginning.'' This suited many
former pagans, for it resembled the supreme God and lesser beings of their old religion and made
Christ an intelligible divine hero; but if the Son's nature differs from the Father's, there are at
least two gods. Arius was an eloquent preacher with a gift for publicity, and his jingles were sung by
dockworkers and schoolchildren. His bishop, Alexander, had a synod condemn and excommunicate him about
320. Arius won the support of Eusebius, bishop of Nicomedia, and the dispute became a contest between
the political capital and the intellectual capital of the East; when he returned, there were riots.

In 325 Constantine summoned a council to Nicaea in Asia Minor. Some three hundred bishops came, many
still bearing the scars of persecution, now traveling in comfort as the emperor's guests. Constantine,
in the jeweled robes of an eastern monarch, opened the sessions, told the bishops that division in the
church was worse than war, and left the matter to them.

\begin{sidebox}{Rejected Trinitarian Schemes}
The early church treated the Trinity as a mystery to be honored, not a riddle to be solved. Two
``Monarchian'' theories protected the one ruling God at the cost of the Father, Son and Spirit needed
to tell the Christian story. \emph{Modalism} saw the three as roles played by God in succession, like
an actor's parts; it could not explain their interaction in the New Testament (as at Jesus' baptism)
and left the real God hidden behind the masks. \emph{Subordinationism} pictured a supreme God assisted
by lesser deities, the Son and the Spirit. Nicaea rejected it in Arius, though the dispute lasted
another half century.
\end{sidebox}

\ssection{True God of True God}
When supporters of Arius stated his case boldly, the council interrupted in dismay: these bishops came
from churches that worshiped Jesus. The hard work was to write a statement that excluded Arius. The
learned historian Eusebius of Caesarea, a friend of the emperor and a halfhearted supporter of Arius,
offered his own creed; the majority wanted something more exact and took another creed, probably from
Palestine, adding the phrases ``true God of true God, begotten not made, of one substance with the
Father.'' The key term \emph{homoousios}, ``of one substance,'' was probably proposed by Hosius of
Cordova, who had the emperor's ear. So arose the Nicene Creed, still the standard of orthodoxy in the
Roman, Eastern, Anglican and other churches. All but two bishops signed; they and Arius were exiled.
Constantine celebrated with a banquet at which bishops reclined beside the emperor among his guards---to
Eusebius it seemed an image of Christ's kingdom---and he kissed the empty eye socket of the
Egyptian bishop Paphnutius.

But afterward Constantine and his successors intervened again and again, banishing and recalling
bishops as one church party or another gained influence at court; doctrine too often depended on the
emperor's favor. The career of Athanasius shows it. Probably present at Nicaea as secretary to
Alexander, he succeeded him as bishop of Alexandria at thirty-three and for fifty years, while the
outcome was uncertain, defended the creed. He was banished five times, and at his loneliest stood, as
the saying goes, ``Athanasius against the world.''

Soon a moderate group, the Semi-Arians, proposed that the Word is ``like'' (\emph{homoiousios}) the
Father rather than ``the same'' (\emph{homoousios}). Athanasius's party insisted on the latter: if
Christ were not fully God he could not fully save. Gibbon sneered that Christians fought over a
diphthong. Shelley replies, with William Hordern's story of a telegram whose missing comma turned
``No, price too high'' into ``No price too high,'' that a single letter can carry immense meaning. Had
the iota entered the creed, Christianity would have had two gods, a Jesus neither God nor man and an
unapproachable Father---one more pagan religion.

\ssection{The Mystery of the Trinity}
Illustrations from nature---egg, plant, the three states of water---miss the personal character of
the doctrine. Its foundation is God's own action in history: revealing himself to Israel, entering
the world in Jesus, coming as the Spirit at Pentecost. Personal analogies are of two kinds. The
\emph{social} analogy starts from three persons so closely bound that they share one life; it is clear
about threeness and must account for unity. The \emph{psychological} analogy starts from one person
with distinct functions; it is clear about unity and vague about threeness. Both were used in the
early church, as by Leonard Hodgson and Karl Barth today.

Between Nicaea and the second general council in 381 the three Cappadocian Fathers---Basil the
Great, Gregory of Nazianzus and Gregory of Nyssa---clarified the term ``person.'' They used the
social analogy but placed the distinctions solely in the inner relations of the one divine Being:
there are not three gods. \emph{Persona} originally meant an actor's mask and did not mean an
individual in today's sense; Christians used it for the living communion of Father, Son and Spirit.
A little later Augustine of Hippo used the psychological analogy: since man is made in the image of
the Trinity, God may be compared to memory, intelligence and will in one mind---only an analogy,
as he knew. Athanasius was not alone after all. He died peacefully at seventy-five, secure in his see
and sure that the creed of Nicaea was the creed of the church.

\begin{sidebox}{Why Is the Trinity Important?}
Father, Son and Spirit are not three names among many but necessary to the Christian story.
\emph{Salvation} meant more than going to heaven: it meant being brought into communion with God, as a
grandson squeezes onto the porch swing between his grandparents and is taken into their love; only if
all three are divine can they bring us into that fellowship. \emph{Worship}: Christians are baptized
into the one name of Father, Son and Spirit, and Paul gives a Trinitarian blessing. \emph{Human
dignity}: against the ancient idea that whoever serves is inferior, the Son serves the Father out of
love, so that service is a mark of Christlikeness, not of weakness.
\end{sidebox}

\begin{keyfacts}
\keypeople{Arius; Alexander of Alexandria; Eusebius of Nicomedia; Constantine; Eusebius of Caesarea;
Hosius of Cordova; Paphnutius; Athanasius; Basil the Great, Gregory of Nazianzus and Gregory of Nyssa
(the Cappadocians); Augustine.}
\keydates{%
  \kd{c. 318}{Arius begins to teach that the Son is a created being}
  \kd{c. 320}{Synod of Alexandria excommunicates Arius}
  \kd{325}{Council of Nicaea; the Nicene Creed (``of one substance with the Father'')}
  \kd{381}{Second general council; the Nicene faith prevails}}
\keyterms{Trinity; Arianism; \emph{homoousios} and \emph{homoiousios}; Semi-Arians; modalism;
subordinationism; Monarchians; person (\emph{persona}); social and psychological analogies; Nicene
Creed.}
\end{keyfacts}
